% === Begin Label Data ===
% ID: 178
% 难度星级: 
% 标签: 
% 备注: 
% 组卷引用次数: 0
% === End  Label Data ===

\begin{problem}{2009}{G}{湖南卷（理）}{21}{数列}
对于数列 $\{u_n\}$, 若存在常数 $M>0$, 对任意的 $n \in \mathbf{N}^*$, 恒有
$$
|u_{n+1}-u_n| + |u_n-u_{n-1}| + \cdots + |u_2-u_1| \le M
$$
则称数列 $\{u_n\}$ 为 $B-$ 数列

（1）首项为 1, 公比为 $q$ $(|q|<1)$ 的等比数列是否为 $B-$ 数列? 请说明理由

（2）设 $S_n$ 是数列 $\{x_n\}$ 的前 $n$ 项和, 给出下列两组论断：

A 组：\circled{1}数列 $\{x_n\}$ 是 $B-$ 数列 \hspace{2em} \circled{2}数列 $\{x_n\}$ 不是 $B-$ 数列

B 组：\circled{3}数列 $\{S_n\}$ 是 $B-$ 数列 \hspace{2em} \circled{4}数列 $\{S_n\}$ 不是 $B-$ 数列

请以其中一组的一个论断为条件, 另一组中的一个论断为结论组成一个命题, 判断所给命题的真假, 并证明你的结论

（3）若数列 $\{a_n\}, \{b_n\}$ 都是 $B-$ 数列, 证明：数列 $\{a_nb_n\}$ 也是 $B-$ 数列
\end{problem}